Let us now address the problem of classification of quartic curves.
For a fixed field $\bbF_q$, $q$ odd, 
we can exploit the previously classified list 
of cubic surfaces as a start. 
For each surface, and for each orbit on points not on lines, 
we perform the projection along the tangent cone to create one quartic curve. 
This guarantees that each isomorphism type of quartic curve with 28 bitangents will be created.
Afterwards, we apply isomorphism testing based on canonical forms. 
This is based on subtructures 
and using canonical froms for graphs using the built-in package Nauty.

\bigskip

Let us look at some examples of this algorithm.
We try $q=7$ and then $q=13.$
We will set  a 
makefile variable to the number of 
(isomorphism types of) cubic surfaces 
with 27 lines over $\bbF_q$. 
For $q=7$, there is exactly one isomorphism type, so we put 

\sourcecodevariable{NB_CUBIC_SURFACES_Q7}

The next command creates a list of 
quartic curves using a projection construction. 
For each isomorphism type of cubic surface, 
and for each point not on any line, we 
consider the intersection of the tangent cone 
with the surface and project onto a plane not containing the point. 
Because of symmetry, it suffices to perform this projection 
only for a set of representatives of the orbits 
on points not on lines. 
Here is the Orbiter command for $q=7$:

\sourcecode{quartic_curves_q7}

The resulting curves are written to file. 
Unfortunately, in this example, 
there is no point which does not lie on any line of the surface.
This means that no quartic curve with 28 lines exists over $\bbF_7.$ 


\bigskip


We move on to the next example, 
which is $q=13$.
Again, we use a makefile variable 
to set the number of isomorphism types of 
cubic surfaces with 27 lines over $\bbF_{13}$. 
There are exactly 4 types:

\sourcecodevariable{NB_CUBIC_SURFACES_Q13}

Just like before, we create all 
quartic curves arising from projections:

\sourcecode{quartic_curves_q13}

Every cubic surface creates a number of quartic curves.
Namely, each orbit of the automorphism 
group on points not on lines of the cubic surface 
creates one quartic curve. 
The curves arising this way are written to file. 
There is one file for each surface.
In Section~\ref{sec:canonical:form:objects:PG}, 
we will continue with the classification of quartic curves.





\bigskip



